\documentclass[book.tex]{subfiles}
\begin{document}

\chapter{Effective Field Theory of Neural Networks}
\label{chap:hopfield}
\noindent\rule{11cm}{0.4pt} \\
\textbf{Prerequisites:}  \autoref{chap:eft},   \\
\textbf{Difficulty Level:} ** \\
\noindent\rule{11cm}{0.4pt}
\vspace{5mm}

We have seen in earlier chapters that an infinitely wide neural network converges to a Gaussian process. This result is of considerable theoretical value, but does not offer much practical guidance to the theory of finite neural networks. 

However, there is an interesting connection to field theory. A classical free field theory is in fact a Gaussian process. There exist a rich variety of theories used in classical physics to model and compute with interactions on top of the free field theory. These same theoretical methods can be used to add corrections from the infinite width Gaussian process network to more practical finite width networks.

%"""
%Large neural networks with many many model parameters work extremely well in practice and thus form the foundation of the modern approach to artificial intelligence. The success of these overparameterized models with far more parameters than training data has led many to simply conjecture that ``more is better'' when it comes to deep learning. However, the practical success of overparameterized models in deep learning appears to be in tension with orthodox machine learning and classic statistical theory. To address this puzzle, in this talk I will examine the notion of model complexity of neural networks -- both at initialization and after training. In particular, looking at the statistics of network outputs, an effective theory approach offers a resolution of this puzzle.
%"""

%References \cite{roberts2021principles}.

\end{document}